% Transcription — fonds Alexandre Grothendieck, shelfmark 109, batch 1 (pages 1-20).
% Established from the facsimile archives/batches/109/batch-01.pdf (a mirror of
% https://grothendieck.umontpellier.fr/109.pdf), per the skill
% transcribe-grothendieck. The batch file carries no cover sheet: PDF page k
% is folder page k. The folder is thirty-seven pages; this is the first of two
% batches.
%
% Pass: Opus 5.5 (claude-opus-5-5), 2026-09-24 — first pass, unchecked against the pages by a human.
%
% Skipped pages: 11, 15 and 19, blank tan separator sheets. The gaps in the
% numbering are the record.
%
% In his hand, transcribed in full: pages 1-8. Seven of them are short
% working sheets preparing the edition: a list of notations with page
% numbers (1), notes for the introduction (2), instructions to the printer
% on symbols (3), an « Index des notations » on a punched card (4), draft
% titles for §§ 130-133 written over a 1982 computer listing (5), a plan of
% the volume with two chapter-dependency diagrams (6), a list of questions
% with page and section numbers (7), and an English « Terminology » list
% with page numbers (8). His lists are given as lists, his page numbers kept.
%
% Pages 9-10: a sheet in another, rounded hand (the same that writes the
% correction list of pp. 16-20 and batch 2), addressed to him (« your
% original »): a proposed English version of his presentation text for the
% « Réflexions Mathématiques » series and « Pursuing Stacks », with the words
% changed « by Breen or by us » underlined in red. It is transcribed in
% full, as a letter to him on the mathematics whose body is his own text in
% revised wording; the red underlines are given as \emph and said once in a
% note. This choice is worth a human look: it could instead be summarised.
% The writer is not named.
%
% Summarised as another's text: pages 12-14 (a three-page typescript
% « COMMENTS ON "PURSUING STACKS" » by a reader who refers to « my paper with
% Heath », English-usage fixes plus substantive remarks on cubical vs
% simplicial sets, Dold-Kan, Loday, Brown-Loday, the nerve; the handwritten
% « so far, so good! » at the foot of p. 14 is in a hand that does not look
% like his and is given inside the note, not as a \marginal), and pages 16,
% 17, 18 and 20 (the English correction list, per the coordinator's ruling
% shared with batch 2: one French \note per page, English usage only in
% aggregate, substantive queries named). Page 20 is a copy of an earlier,
% diagonally cancelled draft of the entries for typescript pp. 87-100,
% dated « Sent 19/12/83 », continued below a rule through p. 115. No mark
% in his hand was found on pages 9-10, 12-14, 16-18 or 20.
%
% What the batch is about. The preparation of the edition of « Pursuing
% Stacks »: the apparatus (dedication, contents, introduction, notes and
% glossary, terminological index, index of notations, bibliography « (?) »),
% the notations (Pi_n(X), (Cat), the double-struck Delta, square and circle
% for the simplex, cube and « globe » categories, A^, K |-> |K|, (Hot),
% Hotab, D(A), K(A), Ch(A)), a terminology list running through primary /
% secondary structure, standard type, n-stacks, coherator,
% Hot-equivalence, test category, modelizer, strict test category,
% aspheric and locally aspheric topos, homotopy interval, elementary and
% local modelizers, each with a page (or section) number; draft titles for
% §§ 130-133 (non-connected bundles, the spherical functor S~, a « crazy
% tentative wrong-quadrant (bi)complex » for the homotopy groups of a
% sphere); a plan « Chap I, Appendice : Lettres à Breen, Chap II à VIII »
% with two Leitfaden; then the correspondents' side — a revised blurb, a
% reader's comments, and the start of the English correction list.
%
% Readings to check first: page 1, « (Ss-sets) » after (Cat) and the sign
% read as \overset{i}{\underset{\ell}{*}}; page 2, « numérotées »; page 3,
% the second and third signs to the printer; page 5, the palimpsest of
% § 131 and the name read « Saleyman » in § 133; page 6, the struck numeral
% between V and VI in the left diagram, and whether the loop means the
% Appendix comes after Chap I; page 8, which line the numbers 13 (second)
% and 25 belong to.

\documentclass[11pt,a4paper]{article}
% Le préambule commun à toutes les transcriptions du fonds Grothendieck.
%
% Il tient en une idée : **l'incertitude se compose, elle ne se résout pas.**
% Une transcription qui lisse un mot douteux a détruit la seule chose pour
% laquelle on la faisait — un lecteur doit pouvoir savoir, à chaque endroit,
% ce qui était sur le feuillet et ce qui a été ajouté. Les six macros qui
% suivent sont donc l'appareil critique, et non de la décoration.
%
% Les mêmes commandes sont reconnues par `scripts/render.mjs`, qui produit la
% vue de lecture du site. Une macro ajoutée ici doit l'être aussi là-bas,
% faute de quoi le rendu échouera bruyamment — ce qui est voulu : un rendu
% silencieusement faux est pire qu'un rendu absent.

\ProvidesPackage{grothendieck}[2026/08/08 Transcriptions du fonds Grothendieck]

\usepackage{amsmath,amssymb,amsthm}
\usepackage{mathrsfs}
\usepackage{stmaryrd}
% Les diagrammes commutatifs sont partout dans ce fonds : c'est la langue
% même de la géométrie algébrique de Grothendieck, et une transcription qui
% les rendrait en prose ne transcrirait rien.
\usepackage{tikz-cd}
% Français par défaut — c'est la langue des feuillets — mais l'anglais est
% chargé aussi : la traduction et le résumé sont des documents anglais, et la
% typographie française (espaces avant les deux-points) y serait une faute.
% Ces documents basculent par \selectlanguage{english} en tête de corps.
\usepackage[english,french]{babel}
\usepackage{fontspec}
\usepackage{geometry}
\geometry{a4paper,margin=25mm,marginparwidth=28mm}
\usepackage{marginnote}
\usepackage{xcolor}
% ulem, not soul. `soul`'s \st and \ul are built on a character-by-character
% scan that XeLaTeX's font machinery breaks: the document fails with an
% undefined control sequence rather than degrading. ulem does the same two jobs
% and is Unicode-engine safe. `normalem` stops it hijacking \emph, which the
% transcriptions use for Grothendieck's own emphasis.
\usepackage[normalem]{ulem}
\usepackage{hyperref}
\hypersetup{colorlinks=true,linkcolor=black,urlcolor=black,pdfborder={0 0 0}}

% --- Métadonnées du lot -----------------------------------------------------
% Reprises telles quelles par le script de rendu, qui les affiche en tête de
% la vue de lecture. Elles fixent ce que le fichier prétend couvrir : sans
% elles, une transcription flotte sans point d'attache dans un fonds de seize
% mille feuillets.

\newcommand{\folder}[1]{\gdef\grothFolder{#1}}
\newcommand{\batch}[1]{\gdef\grothBatch{#1}}
\newcommand{\leaves}[2]{\gdef\grothFirst{#1}\gdef\grothLast{#2}}
\newcommand{\foldertitle}[1]{\gdef\grothFoldertitle{#1}}
\gdef\grothFolder{?}\gdef\grothBatch{?}\gdef\grothFirst{?}\gdef\grothLast{?}\gdef\grothFoldertitle{}

% --- Le résumé --------------------------------------------------------------
%
% Il ouvre la lecture modernisée, sous le titre « Résumé », et il est composé
% en petit. Ce n'est pas un ornement : c'est le seul passage du document qui
% ne soit pas des mathématiques, et un lecteur doit pouvoir le reconnaître
% avant de l'avoir lu — puis le sauter s'il connaît déjà le sujet, ou s'y
% arrêter s'il ne le connaît pas. Le filet qui le suit marque la reprise du
% corps.

\newenvironment{resume}
  {\small\setlength{\parskip}{0.45em}}
  {\par\medskip\hrule\medskip}

% --- Mots-clés ---------------------------------------------------------------
%
% En anglais, en clôture du résumé de la lecture modernisée : le vocabulaire
% moderne sous lequel chercher ce que les feuillets construisent. Le manifeste
% du site (`npm run manifest`) les extrait pour étiqueter la cote entière —
% cette ligne est la source unique des tags, il n'y a pas de fichier à côté.
\newcommand{\keywords}[1]{\par\smallskip\noindent
  {\footnotesize\textsc{Keywords} --- \textit{#1}}\par}

% --- L'appareil critique ----------------------------------------------------

% Le numéro de feuillet, en marge. C'est l'ancre qui permet de régler un
% désaccord sur une formule en regardant *un* feuillet plutôt qu'une chemise
% de six cents. Le site s'en sert aussi pour tourner les pages du fac-similé
% quand on fait défiler la transcription.
\newcommand{\leaf}[1]{%
  \marginnote{\footnotesize\color{gray!70}\textsf{#1}}%
  \label{leaf:#1}\ignorespaces}

% Illisible : on ne devine pas. Un mot inventé qui se lit comme les autres est
% le pire résultat possible.
\newcommand{\ill}{\textcolor{red!60!black}{[\,\ldots\,]}}

% Lecture douteuse : le mot est donné, et signalé comme incertain.
\newcommand{\uncertain}[1]{\uline{#1}}

% Ajout de l'éditeur — une lettre manquante, un mot sous-entendu.
\newcommand{\add}[1]{\textcolor{blue!50!black}{[#1]}}

% Biffé par Grothendieck lui-même. Ce qu'il a rayé fait partie du document :
% une rature dit souvent où la pensée a changé de direction.
\newcommand{\struck}[1]{\sout{#1}}

% Note du transcripteur, sur l'état matériel du feuillet.
\newcommand{\note}[1]{\par\smallskip{\small\color{gray!60!black}\textit{[#1]}}\par\smallskip}

% Note marginale de Grothendieck, distincte du corps du texte.
\newcommand{\marginal}[1]{\par\smallskip\noindent\hspace{1em}%
  {\small\color{gray!60!black}#1}\par\smallskip}

% --- Titre ------------------------------------------------------------------

\AtBeginDocument{%
  \begin{center}
    {\large\bfseries Fonds Alexandre Grothendieck --- cote n\textsuperscript{o}~\grothFolder}\\[2pt]
    {\itshape\grothFoldertitle}\\[4pt]
    {\small feuillets \grothFirst--\grothLast\ (lot \grothBatch)}\\[2pt]
    {\footnotesize\color{gray!60!black}Transcription établie d'après le fac-similé numérisé
     par l'Université de Montpellier.\\
     Les lectures incertaines sont soulignées, les illisibles notées [\,\ldots\,],
     les ajouts entre crochets.}
  \end{center}
  \vspace{1em}}

\endinput


\folder{109}
\batch{1}
\pages{1}{20}
\dating{1983-1984}
\watermark{Édition de démonstration}
\foldertitle{Tables [Corrections et préparation à l'édition] : notes manuscrites (s.d.), copies de notes manuscrites annotées (s.d.), copies de tapuscrit annoté (s.d.), copies de notes manuscrites (1983-1984)}

\begin{document}

\section{Notations}

\page{1}
\note{le titre « Notations » est de sa main, souligné, en haut à droite. Les
signes $\Delta$, $\square$ et $\bigcirc$ de ces feuillets sont tracés d'un
trait doublé sur un côté, à la manière d'une lettre ajourée ; on les rend
par les signes simples. Les nombres à droite sont des renvois de pages}

\begin{itemize}
  \item $\Pi_n(X)$ \quad ? \struck{\ill{}}
  \item $\overset{i}{\underset{\ell}{*}}$, $\check{u}$ \quad ?
    \struck{\ill{}}\note{devant l'astérisque, un petit signe qu'on ne lit
    pas}
  \item $(\mathrm{Cat})$ \uncertain{(Ss-sets)} \quad 16

    \struck{\ill{}} \quad 16 \qquad $\Delta$ \quad $\square$ \quad
    $\bigcirc$
  \item $A_{/X}$, $A^{\circ}$, $A^{\wedge}$ ou $\hat{A}$
  \item $K \mapsto |K|$, $(\mathrm{Ord})$, $(\mathrm{Preord})$,
    \struck{\ill{}} $(\mathrm{Hot})$, $\Delta_n$, $\Delta$, $\Delta[N]$ \quad 22\note{un trait vertical réunit ces six lignes et les renvoie
    au nombre 22}
  \item $\square$, $\square_n$ \quad 29
  \item $\tilde{\square}$, $\tilde{\Delta}$ \quad 34
\end{itemize}

\section{Introduction}

\page{2}
\note{titre de sa main, souligné}

chiffres soulignés dans table des matières. / signe […]

Modifications \uncertain{numérotées} dans lettre à …

Stacks = non commutative chain complexes (cf.\ p.\ 4)

\section{Imprimeur}

\page{3}
\note{titre de sa main, souligné ; en dessous, à gauche, des signes pour
l'imprimeur}

\begin{itemize}
  \item $\underline{N} \longmapsto \mathbb{N}$
  \item $\ell \longmapsto \ell$\note{à gauche, un $\ell$ manuscrit barré
    d'un trait oblique ; à droite, un $\ell$ simple}
  \item \note{un signe fait de deux jambages verticaux accolés, prolongés
    d'un trait horizontal en haut à droite et d'un trait horizontal en bas
    à gauche ; aucune correspondance n'est écrite}
  \item \note{un petit cercle entièrement hachuré ; aucune correspondance
    n'est écrite}
\end{itemize}

\section{Index des notations}

\page{4}
\note{titre de sa main, en haut d'une carte perforée}

\begin{itemize}
  \item $\xrightarrow{\;\approx\;}$ \quad $\xrightarrow{\;\simeq\;}$
  \item $\approx$ \quad $\simeq$
  \item $\underline{\underline{\mathbb{Z}}}$\note{cerclé} \quad
    $\mathbb{R}$ \quad $\mathbb{Q}$ \quad $\mathbb{C}$
  \item $\Delta$ \quad $\square$ \quad $\bigcirc$ \qquad $\Delta_n$ \quad
    $\square_n$ \quad $\bigcirc_n$
  \item $A^{\wedge}$, $A^{\&}$ \qquad $\underline{\mathrm{Hom}}(A,B)$
    $\left\{\begin{array}{l} A, B \in \mathrm{Ob}(\mathrm{Cat}) \\ A, B \in \mathrm{Ob}\, C \text{ catégorie quelc.} \end{array}\right.$
  \item $X_{*}$ \quad $X^{*}$ \quad $L_{\bullet}$ \quad $L^{\bullet}$
  \item $V^{\vee}$
  \item $\mathrm{Hot}$ \quad $\mathrm{Hotab}$\note{ces deux mots d'un
    crayon plus pâle} \qquad $D_{\cdot}(A)$ \quad $D^{\cdot}(A)$ /
    $D^{+}(A)$, $D^{-}(A)$ \struck{\ill{}} $D(A)$

    $\mathrm{Ch}_{\cdot}(A)$, $\mathrm{Ch}^{\cdot}(A)$ \qquad $K^{+}(A)$ \quad $K^{-}(A)$ \quad $K(A)$

    $K_{\cdot}(A)$ \quad $K^{\cdot}(A)$\note{un trait courbe relie
    $\mathrm{Ch}_{\cdot}(A)$, $\mathrm{Ch}^{\cdot}(A)$ à cette dernière
    ligne}
  \item fibering, \struck{\ill{}} cofibering, bifibering functor.
  \item localization (two meanings), colocalization.
\end{itemize}

\section{§§ 130 à 133}

\page{5}
\note{écrit à la main sur une liste d'ordinateur, entre les lignes
imprimées ; le texte imprimé n'est pas reproduit. Ce sont des titres de
paragraphes numérotés, en anglais}

\begin{enumerate}
  \item[130.] A case for non-connected bundles (as components for
    \uncertain{semisimplicial} models).
  \item[131.] \struck{Rambling \ill{} on the} Tentative construction
    of the spherical functor $\tilde{S}$
    \struck{and extension} and extension of notion of « bundles » and
    « multi\uncertain{dimensional} » « extension » of the basic
    \struck{bundle notion}.\note{palimpseste : « Tentative construction
    of the » est écrit au-dessus de « Rambling … », biffé ; la ligne
    « and extension of » est reprise deux fois ; « notion of "bundles" »
    est ajouté en interligne à droite. L'ordre donné ici est celui qu'on
    croit lire, non une certitude. Le numéro 131 est cerclé}
  \item[\struck{132.}] \struck{A crazy formula for homotopy groups
    \ill{} quadrants --- and/or crazy of spheres --- and puzzling about
    hypothetical formula expressing the formula for the homotopy groups
    of spheres.}\note{tout ce bloc est annulé par un grand trait en V et
    ses lignes sont en outre biffées une à une ; la lecture en est
    fragmentaire}
  \item[132.] A crazy tentative wrong-quadrant (bi)complex for the
    homotopy groups of a sphere.\note{numéro cerclé}
  \item[133.] Birth of \uncertain{Saleyman}.\note{la deuxième lettre du
    nom peut être un \emph{u}}
\end{enumerate}

\section{Plan}

\page{6}

\begin{itemize}
  \item Dédicace
  \item Sommaire
  \item Table des matières détaillée
  \item Introduction I

    II\note{une accolade à gauche réunit les deux
    introductions ; la seconde n'a qu'un trait horizontal au lieu du mot}
  \item Chap I \struck{\uncertain{à VIII}}
  \item Appendice : Lettres à Breen
  \item Chap II à VIII\note{« Appendice : Lettres à Breen » est entouré
    d'une boucle qui part de « Chap I » et se ferme sur « Chap II à
    VIII », écrit à droite : l'appendice s'insère, semble-t-il, entre le
    chapitre I et les chapitres II à VIII}
  \item \emph{Notes} et glossaire
  \item Index (terminologique)
  \item Index des notations
  \item Bibliographie (?)
\end{itemize}

\note{au bas de la page, deux schémas de dépendance des chapitres. À
gauche, avec des flèches doubles ; entre V et VI, un chiffre romain
lourdement biffé}

\begin{tikzcd}[column sep=small, row sep=small]
\mathrm{I} & \mathrm{II} \arrow[d, Rightarrow] & & \\
& \mathrm{III} \arrow[d, Rightarrow] & & \\
& \mathrm{IV} \arrow[dr, Rightarrow] \arrow[dd, Rightarrow, bend right=40] & \mathrm{V} \arrow[d, Rightarrow] & \mathrm{VI} \\
& & \mathrm{VII} & \\
& \mathrm{VIII} & &
\end{tikzcd}

\note{à droite, par de simples traits ; VI n'est relié à rien}

\begin{tikzcd}[column sep=small, row sep=small]
& \mathrm{I} \arrow[d, no head] & \\
& \mathrm{II} \arrow[d, no head] & \\
& \mathrm{III} \arrow[d, no head] & \\
& \mathrm{IV} \arrow[dd, no head] \arrow[dddl, no head] & \mathrm{V} \arrow[ddl, no head] \\
& & \mathrm{VI} \\
& \mathrm{VII} \arrow[dd, no head] & \\
\mathrm{VIII} \arrow[dr, no head] & & \\
& \mathrm{IX} &
\end{tikzcd}

\section{Questions}

\page{7}
\note{titre de sa main, souligné ; « p. » renvoie aux pages, « S. » aux
paragraphes}

\begin{itemize}
  \item p.\ 54 \quad S.\ 35 \quad [\uncertain{three} questions on that
    page!]\note{le premier mot est souligné}
  \item p.\ 55 \quad S.\ 35 \quad $x \mapsto \mathrm{Top}(B)$
  \item p.\ 72 \quad S.\ 41 \quad $A^{\wedge}$ cl.\ mod.\ cat ?
  \item p.\ 73, \uncertain{74} / S.\ 42 \quad (Ord) Modelizer ?
\end{itemize}

\section{Terminology}

\page{8}
\note{titre de sa main, souligné, en haut à droite. Les nombres sont dans
une colonne à droite ; leur rattachement aux lignes est donné au plus
près}

\begin{itemize}
  \item primitive, primary, secondary, ternary… structure \quad 4.
    \note{« primitive, » ajouté au-dessus de « primary »}
  \item standard type (iterated fiber products of) \quad 5.
  \item $n$-stacks, $\infty$-stacks, \struck{\ill{}}
    $n$-$\underline{C}$-stacks \quad 13.
  \item coherator \quad 13 \qquad 25
  \item Hot-equivalence of site structure \quad 26, \emph{39}
  \item test category (!) / test-topology \quad 27\note{« category (!) »
    est écrit au-dessus de « topology », et « test » ajouté au-dessus de
    lui}
  \item modelizer ! / model-preserving (or modelizing)
    \struck{\ill{}} functor ! \quad 28 $\to$ \emph{41}, $\to$ \emph{42}
  \item basic modelizer (Cat) / elementary modelizer (!) \quad 29*
  \item \emph{strict} test category \struck{\ill{}} \quad 33, 39
  \item standard simplices / cubes \quad 34
  \item \struck{$\square$, $\square$} \quad 34\note{deux carrés biffés}
  \item aspheric (topos, object of topos, map of objects) \quad 35
  \item locally aspheric topos, tot.\ asph.\ topos
  \item weak equivalence ($f : X \to Y$ map of topos) \quad ?
  \item homotopy interval \quad 37
  \item $[$strict$]$ elementary modelizer \quad 39
  \item local test category / local el.\ modelizer \quad 39
  \item $[$locally$]$ modelizing topos \quad 39
\end{itemize}

\section{Réflexions Mathématiques}

\page{9}
\note{les pages 9 et 10 sont d'une autre main que la sienne, celle qui
écrit la liste de corrections des pages 16 à 20. Elles lui sont adressées
et donnent une version anglaise revue de son texte de présentation. Les
mots soulignés en rouge par le correspondant, qui marquent ce qui a été
changé dans son original, sont rendus en italique ; un petit trait rouge
suit aussi « started » à la page 10. Les mots qu'il écrit en interligne
(« for me », « possible », « thus », « two ») sont remis à leur place}

Reflexions Mathématique\supplied{s}

(The simplest thing for me is to write
out \struck{an} a possible alternative
version as our photocopy of the original is not that clear now. I'll
underline in red places where the words are changed from your original
whether by Breen or by us.)

After a twelve years' silence, the time has come for the author's works of
maturity, with his vision and style renewed. Here is \emph{a} day-by-day
account of an explorer's \struck{travel} journey with occasional
reflections on the travel \emph{itself} (as well as on the traveler and
the \emph{manifold} world around \emph{him,}), thus recapturing its genuine nature: \emph{that} of an impassioned
adventure, rooted in life

In the two forthcoming volumes, under the common title « Pursuing
Stacks » (the first volume of which is subtitled « The Modelizing
Story, »), the author sketches some themes of a vast sythesis\note{sic} of
homotopical algebra, (co)homological algebra and topos theory --- ripe now
for more than fifteen years \emph{but} never yet begun. The leading
thread, tenacious and omnipresent (while often remaining implicit) comes
from algebraic geometry, and from the intuitions (rich precise and
\emph{provocatively} fragmentary) \emph{surrounding} the notions of
sheaves, stacks, 2-stacks « etc. » on a topos.

Thes\supplied{e} two volumes are the first in a planned, considerably wider
series of « Réflexions Mathématiques » \emph{in which} the author intends
to \emph{present}, among other things, some « naive ideas » about two
« new continents », eager to be discovered and explored. One \emph{stems
from}

\page{10}
a suitable notion of a « tame » topological space, viewed as the natural
set-up for expressing and stimulating the topological intuition of shape
and form, which for so long has been squeezed \emph{into} the
straitjacket of a « General Topology » tailored for the use of the
analyst, not for geometry. This is inapt notably (because of pervasive
« wildness » phenomena) for \emph{the study}\struck{ing} \emph{of}
homotopy types of spaces of homeomorphisms, imbeddings \& the like, or for
developing a comprehensive « dévissage » theory of stratified structures
(whereas such structures \struck{abound} \emph{crop up everywhere} notably
in any question of « moduli » \struck{of} for geometrical configurations
capable of continuous variation and giving rise to degeneracy
phenomena). The other continent appears \emph{at} the confluence of a
multitude of streams of thought and insights, coming from topology,
conformal geometry, algebraic geometry, discrete and profinite groups,
algebraic groups over number fields, regular polyhedra, arithmetics,
« motives »…, \emph{called together for} the exploration of the action
of profinite « absolute » Galois groups (such as
$\mathrm{Gal}_{\bar{\mathbb{Q}}/\mathbb{Q}}$) on certain (so called
« anabelian ») profinite geometrical fundamental groups (such as
$\pi_1(\mathbb{P}_{\bar{\mathbb{Q}}} - \{0,1,\infty\})$). One of the
leading threads comes from the combinatorial topology

(Continues without change until last paragraph)

The first \struck{two} two volumes « Pursuing Stacks », started through the impetus of a correspondence in English are
written in that language. They begin with a substantial common
introduction \emph{of which a first part}, of personal rather than
technical character, is written in French (the author's adopted
language.).

\section{Comments on « Pursuing Stacks »}

\page{12}
\note{les pages 12 à 14 sont un tapuscrit de trois pages, intitulé
« COMMENTS ON "PURSUING STACKS" », d'un lecteur qui renvoie à « my paper
with Heath » et n'est pas nommé ici. Il est résumé page à page. On n'y a
relevé aucune marque de sa main}

\note{d'abord une liste de mots à corriger, repérés par page et ligne du
tapuscrit (pp.\ 2 à 4 de l'introduction, puis « Section 14 etc », pp.\ 2
et 3) : pour l'essentiel des fautes d'orthographe anglaise
(\emph{inadequate}, \emph{pertinent}, \emph{misled}, \emph{CW-spaces},
\emph{adequate}…). Puis une remarque sur la p.\ 3, « cubical versus
simplicial sets » : les ensembles simpliciaux doivent leur popularité à
ce que leur catégorie est cartésiennement fermée, que la réalisation
donne une équivalence entre la catégorie homotopique des complexes de Kan
et celle des CW-espaces, qu'elle commute au produit (pour la topologie
compactement engendrée), à l'existence d'un théorème de Dold--Kan et
d'une démonstration simpliciale de la suite spectrale de Serre ; les
ensembles cubiques, commodes pour certains auteurs (un exposé de
l'homologie singulière, les méthodes de Brown--Higgins), sont moins
établis dans la littérature}

\page{13}
\note{suite : pour les ensembles cubiques, l'équivalence des catégories
homotopiques n'est démontrée en détail que dans un mémoire non publié ;
leur produit cartésien n'a pas la bonne réalisation, seul un produit
tensoriel l'a ; un théorème de Dold--Kan demande la structure
supplémentaire des « connexions » de Brown--Higgins. À la p.\ 4, le
lecteur préférerait parler du théorème de Dold--Kan plutôt que de Puppe,
et rappelle que Kan a décrit le foncteur des complexes de chaînes vers
les groupes abéliens simpliciaux, en dimension $n$, par
$\mathrm{Hom}(C_N \Delta^n, C)$. Il s'étonne que l'approche de Loday des
types d'homotopie passe par les ensembles $n$-simpliciaux, et note que
la théorie de van Kampen de Brown--Loday utilise les groupoïdes
$n$-uples, plus précisément les $\mathrm{Cat}^n$-groupes ; il se demande
si, pour définir de tels objets dans d'autres situations, comme les
topos, il ne faudrait pas des constructions d'intervalle unité, là où
interviendraient les structures de champs. Suivent des retouches de
style (pp.\ 4 et 5), et, p.\ 7, un renvoi aux travaux publiés sur le nerf
et son rapport aux types d'homotopie ; p.\ 8, le paragraphe central est
jugé intéressant}

\page{14}
\note{suite : trouver des catégories qui modélisent des situations
géométriques puis en prendre le nerf s'est révélé utile en topologie
géométrique (K-théorie algébrique des espaces), sans être encore un outil
de base ni avoir mené à des calculs ; un test proposé : déterminer
$\pi_3(S^2)$ sans connaître la fibration de Hopf $S^1 \to S^3 \to S^2$.
Puis des retouches de mots pour les pp.\ 9 à 17 (\emph{developing},
\emph{standard} pour \emph{standart}, qui revient souvent, \emph{the
formal analogy}, « commute with formation »…), et une question, p.\ 13 :
la catégorie de tous les groupoïdes n'est-elle pas un topos parce
qu'elle n'a pas de classifiant des sous-objets ? Au pied de la page,
d'une main qui ne semble pas la sienne : « so far, so good! »}

\section{Liste de corrections}

\page{16}
\note{les pages 16 à 18 et 20 sont d'une autre main que la sienne, celle
des pages 9 et 10 : le début d'une liste, en anglais, de corrections
page et ligne par page et ligne, qui lui est adressée et se poursuit dans
le lot 2 du dossier. Le tapuscrit visé n'est pas nommé sur les feuillets
(sans doute celui de \emph{Pursuing Stacks}). Elle est résumée page à
page ; on n'y a relevé aucune marque de sa main}

\note{intitulé « Errors in 1st 100 pages of notes + letter to Quillen ».
Pour la lettre à Quillen (pp.\ 5, 5', 6') : « commute with » et non
« commute to », « standard » et non « standart », un article manquant.
Pour les « Reflections » (pp.\ 1 à 41) : des corrections d'usage et
d'orthographe (\emph{indispensable}, \emph{impediments},
\emph{manifold}, \emph{adequate}, \emph{reminds}, « we have » plutôt que
« we got »…) ; p.\ 32, « probably "modelizer" » ; p.\ 34, la phrase
« which allow… » est jugée obscure}

\page{17}
\note{suite, pp.\ 41 à 90 du tapuscrit : « aspheric » et non « asperic »
(le correspondant ajoute en français « avec les asperges ! »),
\emph{analogues}, \emph{denote}, \emph{Assume}, les contractions
(« isn't », « doesn't ») mal placées, « last but not least », « compare
with » comme « commute with » ; p.\ 66, un doute sur le latin
(\emph{compositum}, \emph{composita} ou \emph{compositae})}

\page{18}
\note{suite, pp.\ 91 à 100, feuillet numéroté « -2- » : « in view of »
et « with a view to », « we've got » ou « we have got », « avatars » et
non « avators », « informal » ; p.\ 99, la remarque que $X$ et $Y$ sont
devenus $F$ et $G$ ; p.\ 100, un article manquant}

\page{20}
\note{copie d'un état antérieur des entrées pour les pp.\ 87 à 100,
annulé par deux traits obliques, avec au bas la mention « Sent
19/12/83 » ; puis, sous un trait, la suite pour les pp.\ 102 à 115 :
de nouveau « commute with », les contractions, \emph{developing},
\emph{relationships}, et p.\ 107 « Kan » et non « Khan »}

\end{document}
